% ==STYLE-BEGIN==
\documentclass[UTF8,11pt]{ctexart}
\usepackage{amsmath,amssymb,amsthm}
\usepackage[a4paper,margin=2.4cm,top=2.3cm,bottom=2.3cm]{geometry}
\usepackage{fancyhdr}
\usepackage{enumitem}
\usepackage{hyperref}
\hypersetup{hidelinks}

% ---- 版式 ----
\linespread{1.15}
\setlength{\parskip}{0.35em}
\setlength{\parindent}{1.6em}
\setlist{nosep,leftmargin=2em}

% ---- 页眉 ----
\pagestyle{fancy}
\fancyhf{}
\fancyhead[L]{\small\coursename}
\fancyhead[R]{\small\thepage}
\renewcommand{\headrulewidth}{0.4pt}

% ---- 定理环境：中文用正体，标题加粗，按小节共用一条流水号 ----
\newtheoremstyle{cnthm}{0.7em}{0.7em}{\normalfont}{}{\bfseries}{}{0.5em}{}
\theoremstyle{cnthm}
\newtheorem{thm}{定理}[section]
\newtheorem{lem}[thm]{引理}
\newtheorem{prop}[thm]{命题}
\newtheorem{cor}[thm]{推论}
\newtheorem{defn}[thm]{定义}
\newtheorem{rem}[thm]{注}
\newtheorem{exa}[thm]{例}
\newtheorem{conj}[thm]{猜想}
\renewcommand{\proofname}{\bfseries 证明}
\providecommand{\coursename}{}
% ==STYLE-END==
\renewcommand{\coursename}{数学分析\,·\,0914}
\title{数学分析 2024-09-14}
\usepackage{graphicx}
\newtheorem{theorem}{定理}
\newtheorem{example}{例}
\date{}
\begin{document}
\maketitle

\section{Recall}

% 原文图片：Figure_1.png；图片文件未附
% 若图片已上传到项目，取消下一段注释即可显示
% \begin{figure}[htbp]
%   \centering
%   \includegraphics[width=0.6\textwidth]{Figure_1.png}
%   \caption{原文图片}
% \end{figure}

将曲边梯形的面积看作无限个小长方形面积的和，取 \(x_i=i/N\)，有
\[
\lim_{N\to\infty}
\sum_{i=0}^{N-1} f(x_i)\frac1N
=
\int_0^1 f(x)\,\mathrm dx,
\]
前提是相应积分存在。

\paragraph{问题.}
如果 \(f\) 与 \(g\) 在 \([0,1]\) 上仅有有限个点的函数值不同，是否仍有
\[
\int_0^1 f(x)\,\mathrm dx
=
\int_0^1 g(x)\,\mathrm dx?
\]
在两者均可积的前提下，答案是肯定的。因为改变有限个点的函数值不改变积分的值。

\section{丢掉高阶无穷小}

以 \(f(x)=x^2\) 为例。函数在 \(x=1\) 处的极限为
\[
\lim_{x\to1} f(x)=\lim_{x\to1} x^2=1.
\]
求导数时，
\begin{align*}
f'(1)
&=\lim\limits_{\Delta x\to0}
  \frac{f(1+\Delta x)-f(1)}{\Delta x}\\
&=\lim\limits_{\Delta x\to0}
  \frac{(1+\Delta x)^2-1}{\Delta x}.
\end{align*}
分子可以展开为
\[
f(1+\Delta x)-f(1)
=
2\Delta x+(\Delta x)^2.
\]
因此
\[
f'(1)
=
\lim_{\Delta x\to0}
 \frac{2\Delta x+(\Delta x)^2}{\Delta x}
=
\lim_{\Delta x\to0}(2+\Delta x)=2.
\]
其中 \((\Delta x)^2\) 是相对于 \(\Delta x\) 的高阶无穷小。

另设
\[
g(x)=
\begin{cases}
2,&x=1,\\
x^2,&x\ne1,
\end{cases}
\]
虽然 \(g(1)=2\)，仍有 \(\lim\limits_{x\to1} g(x)=1\)。

\section{有限与无限}

\begin{enumerate}
\item \textbf{鸽笼原理：}将 \(N+1\) 只鸽子放进 \(N\) 个笼子，
必有一个笼子至少装有两只鸽子，其中 \(N\) 为有限正整数。
\item \textbf{希尔伯特旅馆：}若无穷多个房间均已住满，
仍可令原住在第 \(n\) 间的客人搬至第 \(n+1\) 间，
从而腾出第 \(1\) 间。
\end{enumerate}

\section{数学归纳法}

要证明命题 \(P(n)\) 对所有 \(n\in\mathbb Z^+\) 成立，
只须证明：
\begin{enumerate}
\item \(P(1)\) 成立；
\item 对任意 \(n\in\mathbb Z^+\)，若 \(P(n)\) 成立，
则 \(P(n+1)\) 也成立。
\end{enumerate}

\section{数系的扩充}

\section{从自然数到整数}

自然数对加法封闭，但减法未必仍得到自然数，例如
\(1-2=-1\)。为了使加法方程 \(n+x=m\) 总能求解，
需要引入整数：
\[
\mathbb N\longrightarrow\mathbb Z.
\]

\section{从整数到有理数}

整数对加法和乘法封闭，但方程 \(nx=m\) 在整数中
未必有解。例如 \(2x=1\) 不存在整数解。
当 \(n\ne0\) 时，在有理数中可以解这类方程：
\[
\mathbb Z\longrightarrow
\mathbb Q
=
\left\{\frac nm:n,m\in\mathbb Z,\ m\ne0\right\}.
\]

继续扩充数系，可得到
\[
\mathbb Q\longrightarrow\mathbb R\longrightarrow\mathbb C.
\]
实数的引入涉及连续性；复数使每个非常数复系数
多项式在 \(\mathbb C\) 中有根。

\section{布尔巴基学派的结构观}

数学结构可从以下几个方面考察：
\[
\left\{
\begin{array}{l}
\text{代数结构},\\
\text{序结构},\\
\text{拓扑结构}.
\end{array}
\right.
\]

\end{document}

\end{document}
